%\documentclass[handout]{ximera}
\documentclass[nooutcomes,noauthor]{ximera}

%% page layout
\usepackage[in,headings]{fullpage}
\raggedright
\setlength\headheight{13.6pt}


%% fonts
\usepackage{euler}

\usepackage{FiraMono}
\renewcommand\familydefault{\ttdefault} 
\usepackage{mathastext}
\usepackage[htt]{hyphenat}

\usepackage[T1]{fontenc}
\usepackage[scaled=1]{FiraSans}

\usepackage{wedn}
\usepackage[T1]{fontenc}

%% wrap text around scripts
\usepackage{wrapfig}

\tikzset{>=stealth}
%% snap! scripts
\usepackage{scratch3}

\usepackage{adjustbox}

%% journal style
\makeatletter
\newcommand\journalstyle{%
  \def\activitystyle{activity-chapter}
  \def\maketitle{%
    \addtocounter{titlenumber}{1}%
                {\flushleft\small\sffamily\bfseries\@pretitle\par\vspace{-1.5em}}%
                {\flushleft\LARGE\sffamily\bfseries\thetitlenumber\hspace{1em}\@title \par }%
                {\vskip .6em\noindent\textit\theabstract\setcounter{question}{0}\setcounter{sectiontitlenumber}{0}}%
                    \par\vspace{2em}
                    \phantomsection\addcontentsline{toc}{section}{\thetitlenumber\hspace{1em}\textbf{\@title}}%
                     }}
\makeatother



%% thm like environments
\let\question\relax
\let\endquestion\relax

\newtheoremstyle{QuestionStyle}{\topsep}{\topsep}%%% space between body and thm
		{}                      %%% Thm body font
		{}                              %%% Indent amount (empty = no indent)
		{\bfseries}            %%% Thm head font
		{)}                              %%% Punctuation after thm head
		{ }                           %%% Space after thm head
		{\thmnumber{#2}\thmnote{ \bfseries(#3)}}%%% Thm head spec
\theoremstyle{QuestionStyle}
\newtheorem{question}{}



\let\freeResponse\relax
\let\endfreeResponse\relax

%% \newtheoremstyle{ResponseStyle}{\topsep}{\topsep}%%% space between body and thm
%% 		{\wedn\bfseries}                      %%% Thm body font
%% 		{}                              %%% Indent amount (empty = no indent)
%% 		{\wedn\bfseries}            %%% Thm head font
%% 		{}                              %%% Punctuation after thm head
%% 		{3ex}                           %%% Space after thm head
%% 		{\underline{\underline{\thmname{#1}}}}%%% Thm head spec
%% \theoremstyle{ResponseStyle}

\usepackage[tikz]{mdframed}
\mdfdefinestyle{ResponseStyle}{leftmargin=1cm,linecolor=black,roundcorner=5pt,frametitlefont=\wedn\bfseries,%frametitle={\underline{\underline{Response}}:}
, font=\wedn\bfseries,}%\begin{mdframed}[style=mystyle]foo\end{mdframed}


\ifhandout
\NewEnviron{freeResponse}{}
\else
%\newtheorem{freeResponse}{Response:}
\newenvironment{freeResponse}{\begin{mdframed}[style=ResponseStyle]}{\end{mdframed}}
\fi



%% attempting to automate outcomes.

\newwrite\outcomefile
  \immediate\openout\outcomefile=\jobname.oc
\renewcommand{\outcome}[1]{\edef\theoutcomes{\theoutcomes #1~}%
\immediate\write\outcomefile{\unexpanded{\outcome}{#1}}}

%% \newcommand{\outcomelist}{\begin{itemize}\theoutcomes\end{itemize}}



%% my commands

\newcommand{\snap}{{\bfseries\itshape\textsf{Snap!}}}
\newcommand{\flavor}{\link[\snap]{https://snap.berkeley.edu/}}


\usepackage{tkz-euclide}
\tikzstyle geometryDiagrams=[rounded corners=.5pt,ultra thick,color=black]
\colorlet{penColor}{black} % Color of a curve in a plot

\usepackage{fullpage}
\makeatletter
%% no number for activity
\newcommand\logostyle{%
  \def\activitystyle{activity-chapter}
  \def\maketitle{%
                {\flushleft\small\sffamily\bfseries\@pretitle\par\vspace{-1.5em}}%
                {\flushleft\LARGE\sffamily\bfseries\@title \par }%
                {\vskip .6em\noindent\textit\theabstract\setcounter{problem}{0}\setcounter{sectiontitlenumber}{0}}%
                    \par\vspace{2em}
                    \phantomsection\addcontentsline{toc}{section}{\textbf{\@title}}%
                     \setcounter{titlenumber}{0}}}
\makeatother
\newcommand{\nameblankgen}{\noindent\textbf{Name(s) (please print):}\ \hrulefill \\

\hrulefill}
\logostyle



%% %%% Graph paper background
%% \def\mygraphpaper{%
%%   \begin{tikzpicture}
%%     %\draw[line width=.4pt,draw=black!30] (0,0) grid[step=1mm] (\paperwidth,\paperheight);
%%     \draw[line width=.4pt,draw=blue!50] (0,0) grid[step=.2in] (\paperwidth,\paperheight);
%%   \end{tikzpicture}%
%% }
%% \usepackage{background}
%% \backgroundsetup{
%%   angle=0,
%%   contents=\mygraphpaper,
%%   color=black,
%%   scale=1,
%% }





\pagenumbering{gobble}


\title{There can be only seven}
\author{Bart Snapp}

\outcome{Learning outcome goes here.}

\begin{document}
\begin{abstract}
  We classify all frieze patterns.
\end{abstract}
\maketitle

\noindent\textbf{Group members (please print):}\ \hrulefill \\

\hrulefill

The title of this section says it all---there are only seven frieze
groups. Believe it or not, we will tackle this monster by solving an
unrelated problem!

\section{Pascal's Pizza Shop}

\textit{Pascal's Pizza Shop} makes a variety of pizzas, all which come
with cheese.  However, they do not offer double toppings---no ``double
cheese'' and so on.

\begin{problem}
If the shop offers two additional toppings, namely basil and garlic.
\begin{enumerate}
\item How many different 1-item pizzas (cheese plus one other
  topping) can he make?  List the possibilities.
\item How many different 2-item pizzas can he make?  List them.
\item How many different 0-item pizzas can he make?
\item In total, how many different pizzas can he make?
\end{enumerate}
\end{problem}


\begin{problem}
Now suppose the shop has basil, garlic, and spinach toppings
available.
\begin{enumerate}
\item How many 0-item pizzas can he make?
\item How many 1-item pizzas can he make?
\item How many 2-item pizzas can he make?
\item How many 3-item pizzas can he make?
\item In total, how many different pizzas can he make?
\end{enumerate}
\end{problem}


\begin{problem}
Complete the following chart by listing and counting the
possibilities, where $n$ is the number of toppings Pascal has
available, and $k$ is the number of toppings used.\\
\[
{\renewcommand{\arraystretch}{2}
\begin{array}{|c||c|c|c|c|c|c||c|}
    \hline
          & k=0 & k=1 & k=2 & k=3 & k=4 & k=5 & \text{Total}\\
    \hline\hline
    n=0 &       &       &       &       &       &    &  \\
    \hline
    n=1 &       &       &       &       &       &    &  \\
    \hline
    n=2 &       &       &       &       &       &    &  \\
    \hline
    n=3 &       &       &       &       &       &    &  \\
    \hline
    n=4 &       &       &       &       &       &    &  \\
    \hline
    n=5 &       &       &       &       &       &    &  \\
    \hline
\end{array}}
\]
\end{problem}


\begin{problem} 
Often the numbers in the table you filled in above are arranged a bit
differently, as below, where the top ``1'' counts the number of
$0$-topping pizzas that can be made if zero toppings are
available. Complete the first 8 rows of the triangle, which is called
\textit{Pascal's Triangle}.\index{Pascal's Triangle}
\[
\begin{tabular}{c c c c c c c}
      &   &   & 1 &   &   &  \\
      &   & 1 &   & 1 &   &  \\
      & 1 &   & 2 &   & 1 &  \\
    1 &   & 3 &   & 3 &   & 1
\end{tabular}
\vspace{2in}
\]
Do you see some patterns in the triangle?  Describe them and try to
explain why (in terms of pizzas) why they occur.
\end{problem}




\section{Symmetries as toppings}

We are going to count the various way that a frieze pattern could have
symmetry. Let's list the basic geometric transformations that can
preserve frieze patterns:
\begin{enumerate}
\item Horizontal translations denoted by $\mat{T}$.
\item Horizontal reflections denoted by $\mat{F_h}$.
\item Vertical reflections denoted by $\mat{F_v}$.
\item Glide reflections denoted by $\mat{G}$.
\item Rotations of $180^\circ$ denoted by $\mat{R}$.
\end{enumerate}
Hence, just as we counted pizza's with a number of toppings, we can
count groups of symmetries, where we select groups containing the
transformations above.

To start, note that every frieze pattern has symmetry through horizontal
translations. So we don't need to count that one---this is like saying
that every pizza has cheese on it. Let's list out all possible combinations
of the remaining 4 transformations. To make sure we get them all, we are
going to do this in a very orderly fashion.

\begin{problem} 
How many combinations of symmetries should there be with a single
transformation?  What are they?
\end{problem}


\begin{problem} 
How many combinations of symmetries should there be with two
transformations?  What are they?
\end{problem}


\begin{problem} 
How many combinations of symmetries should there be with three
transformations?  What are they?
\end{problem}


\begin{problem} 
How many combinations of symmetries should there be with four
transformations?  What are they?
\end{problem}




\section{Grouping by groups}

Now that you have a listing of all possible combinations of the
symmetries of the frieze patterns, let's literally try to group them
by groups---we're going to put the combinations you already found into
the groups we already know. Here are two leading questions to help you
out:


\begin{problem} 
Why must a frieze pattern have symmetry through glide reflections, if
you know it has symmetry through both translations and vertical
reflections?
\end{problem}

I'll take this one, just to show you how it's done. We'll imagine some
sort of generic frieze pattern, and draw the following pictures:
\[
\includegraphics{friezeTF-G.pdf}
\]
Hence, from the picture we see that if the frieze pattern has symmetry
through both $\mat{T}$ and $\mat{F_v}$, then it must also have
symmetry through $\mat{G}$.


\begin{problem} 
Why must a frieze pattern have symmetry through $180^\circ$ rotations,
if you know it has symmetry through translations, glide reflections,
and horizontal reflections?
\end{problem}



\begin{problem} 
Why must a frieze pattern have symmetry through glide reflections, if
you know it has symmetry through translations, horizontal reflections,
and $180^\circ$ rotations?
\end{problem}


\begin{problem} 
Why must a frieze pattern have symmetry through glide reflections, if
you know it has symmetry through translations, horizontal reflections,
and vertical reflections?
\end{problem}


\begin{problem} 
Why must a frieze pattern have symmetry through $180^\circ$ rotations, if
you know it has symmetry through translations, horizontal reflections,
and vertical reflections?
\end{problem}



\begin{problem} 
Why must a frieze pattern have symmetry through vertical reflections,
if you know it has symmetry through translations, horizontal
reflections, and $180^\circ$ rotations?
\end{problem}





\begin{problem} 
Can you figure out which frieze groups go with which collections of
symmetries you listed above?
\end{problem}



From this we can finally give a rigorous definition of a frieze pattern
as promised:

\begin{definition}\index{frieze pattern}
A \dfn{frieze pattern} is a pattern whose symmetries are given by
one of the seven frieze groups.
\end{definition}




\paragraph{Some Observations and Thoughts}


Now that we have classified all of the frieze groups, let me tell you about
another classification project. Around 1971, a mathematician by the
name of Gorenstein sought to understand \textbf{all} finite
groups. To do this, he proposed that we attempt to classify all of
what are called ``finite simple groups.'' While the definition of a
\textit{simple group} is somewhat beyond the scope of this course, to
put it in SAT terms:
\begin{center}
\textit{simple group} is to \textit{group} \qquad as \qquad \textit{prime number} is to \textit{number}
\end{center}
This goal was eventually finished in 2004---however, the proof of the
classification was enormous. It spanned 15000 pages and required the
work of over 100 mathematicians. Some of the arguments are described
as ``diabolically difficult'' and the proof surely presses to the
limits of human thought.



\end{document}
